\section{综合算例与关键结论}
\label{l15:sec:examples}

\subsection{非整块输入的完整归约过程}

考虑 $N=1000$ 个全部等于 $1$ 的输入，取 $B=128$、$C=2$，使用两阶段线程束归约。每个块处理 $2CB=512$ 个逻辑位置，所以需要 $\lceil1000/512\rceil=2$ 个块。第 $0$ 块处理下标 $0$ 至 $511$，全部位置有效；第 $1$ 块从下标 $512$ 开始，只有 $488$ 个真实输入，其余 $24$ 个逻辑位置补零。

尾块线程 $t$ 依次尝试读取 $512+t$、$640+t$、$768+t$ 和 $896+t$。前三次读取对所有 $0\leq t<128$ 都有效；第四次读取要求 $896+t<1000$，即 $t\leq103$。因此线程 $0$ 至 $103$ 的寄存器部分和为 $4$，线程 $104$ 至 $127$ 的部分和为 $3$。边界处，线程 $103$ 读取到最后一个有效下标 $999$，线程 $104$ 的最后一次读取被边界判断跳过。

前三个线程束分别包含 $32$ 个局部值 $4$，各得到 $128$。最后一个线程束包含线程 $96$ 至 $127$，其中八个局部值为 $4$，二十四个局部值为 $3$，因此得到 $8\times4+24\times3=104$。与完整块相比，两块的中间结果如表\ref{l15:tab:tail-example} 所示。

\begin{table}[htbp]
\centering
\caption{1000个全1输入在两个线程块中的归约结果}
\label{l15:tab:tail-example}
\begin{tabularx}{\linewidth}{@{}lYY@{}}
\toprule
阶段或数量 & 第 $0$ 块 & 第 $1$ 块\\
\midrule
有效输入数 & $512$ & $488$\\
逻辑补零数 & $0$ & $24$\\
四个线程束的和 & $128,128,128,128$ & $128,128,128,104$\\
第一个线程束第二次归约 & $512$ & $488$\\
\bottomrule
\end{tabularx}
\end{table}

在两个块的第二阶段，第一线程束均先读取四个共享部分和，再用 $28$ 个零补齐为 $32$ 项。它们共同产生块和 $512$ 与 $488$。若使用原子输出，两个块分别向已经清零的总和地址执行一次原子加法，最终得到 $1000$。

若使用部分和输出，第一轮的结果数组是 $[512,488]$，第二轮再对这个长度为 $2$ 的数组归约，长度变化为 $1000\to2\to1$。第二轮仍可启动 $128$ 个线程，只有线程 $0$、$1$ 的首次输入读取有效，其余逻辑输入均取 $0$；所有线程继续遵守原有线程束与块屏障结构。第一次线程束归约得到一个 $1000$ 和三个 $0$，第二次线程束归约得到最终的 $1000$。

这个算例说明，输入长度不足一个块或不足一个线程束时，仍可以保持固定的线程组织，只在数据加载处改变有效值与单位元的分配。有效输入完整参与、共享数组初始化以及完整线程束参与由此能够同时成立。

\subsection{大规模输入需要多少轮核函数}

考虑 $N=1{,}000{,}000$，取 $B=256$、$C=4$，每块容量为 $M=2048$。由于 $1{,}000{,}000=488\times2048+576$，第一轮有 $488$ 个完整块和一个尾块，共产生 $489$ 个部分和。尾块的段首为 $999424$，包含 $576$ 个有效输入，其余 $1472$ 个逻辑位置补零。

若继续用同样的配置归约部分和数组，第二轮只需一个块，因为 $489<2048$。完整的长度变化和块数如表\ref{l15:tab:rounds} 所示。

\begin{table}[htbp]
\centering
\caption{百万元素的多轮分块归约}
\label{l15:tab:rounds}
\begin{tabular}{@{}lrrrr@{}}
\toprule
轮次 & 输入长度 & 每块容量 & 启动块数 & 最后一个块的有效输入数\\
\midrule
第一轮 & $1{,}000{,}000$ & $2048$ & $489$ & $576$\\
第二轮 & $489$ & $2048$ & $1$ & $489$\\
\bottomrule
\end{tabular}
\end{table}

若输入也全部为 $1$，第一轮前 $488$ 个块各输出 $2048$，尾块输出 $576$；第二轮将这些值合为 $488\times2048+576=1{,}000{,}000$。这个输入规模在理想二叉树中具有约 $20$ 层依赖，但多轮实现只启动两次核函数，因为每次核函数内部已经完成了每线程局部累加和块内的多级归约。核函数轮次、树的算术深度和块内屏障阶段数分别描述不同层次的工作。

若改用原子输出，同样的第一轮仍需 $489$ 个块，但每块直接贡献一次原子加法，从而在一次核函数中完成全局结果。此时原子更新数量为 $489$，而逐元素原子方案需要 $1{,}000{,}000$ 次。是否减少了一次核函数启动、是否保留最终同地址竞争，应一起计入方案比较。

\Needspace{24\baselineskip}
\subsection{归约结构与开销的联系}

为统一比较，表\ref{l15:tab:summary} 取 $B\geq64$ 且为二的幂，统计每块的同步阶段数；一次块屏障由参与线程共同执行，表中没有把它按线程重复累计。直接树保留前面代码中的末次屏障。两阶段方法本身只要求由完整线程束组成，表中的二的幂条件用于与其他树形实现采用相同配置。

\begin{table}[H]
\centering
\small
\caption{不同块内实现的存储与同步开销}
\label{l15:tab:summary}
\begin{tabularx}{\linewidth}{@{}>{\raggedright\arraybackslash}p{3.1cm}>{\raggedright\arraybackslash}p{2.8cm}>{\raggedright\arraybackslash}p{2.7cm}Y@{}}
\toprule
实现 & 中间存储 & 每块屏障阶段数 & 主要变化\\
\midrule
递增步长全局树 & 全局数组 & $1+\log_2B$ & 活动线程及地址逐层稀疏\\
递减步长全局树 & 全局数组 & $1+\log_2B$ & 连续地址与连续活动前缀\\
共享内存树 & $B$ 个共享值 & $1+\log_2B$ & 中间读写留在共享内存\\
共享树与末段交换 & $B$ 个共享值\newline 及寄存器 & $\log_2B-5$ & 最后64项交给一个线程束\\
两阶段线程束归约 & $B/32$ 个共享值\newline 及寄存器 & $1$ & 各束独立归约，只在交接处同步\\
粗化与两阶段归约 & $B/32$ 个共享值\newline 及寄存器 & $1$ & 每线程先处理 $2C$ 个输入，减少块数\\
\bottomrule
\end{tabularx}
\end{table}

归约优化涉及三个相互联系的选择：怎样把二元合并分配给线程、在哪一级存储中传递部分和，以及让哪些线程共同同步。递减步长改善线程与地址的布局；共享内存和线程束交换逐步降低中间通信代价；线程粗化把更多工作放到单线程内部，以减少重复的块内归约。它们分别作用于不同的开销，因此可以组合使用。

正确性可以沿输入覆盖和阶段依赖追踪：每个有效输入进入一个且仅一个局部和，跨线程读取发生在所需结果已可见之后，每个块结果被全局归约恰好使用一次。改变存放位置、分组和同步范围时，这些关系需要同时保留。浮点输入还应按所要求的误差标准验证数值结果，因为不同的合法归约顺序可能产生不同的舍入误差。
